%%%%%%%%%%%%%%%%%%%%%%%%%%%%%%%%%%%%%%%%%%%%%%%%%%%%%%%%%%%%%%%%%%%%%%
% Lớp 9 | Chương V. Đường tròn | Cung và dây của một đường tròn
% Dạng 9H5?14-2: Góc ở tâm, số đo cung
%%%%%%%%%%%%%%%%%%%%%%%%%%%%%%%%%%%%%%%%%%%%%%%%%%%%%%%%%%%%%%%%%%%%%%


% nguồn: dan_tu_AI (nhập 14/11/2026)
\begin{ex}%[9H5H14-2]%[ma:2846d8d4]
	Cho hai điểm $A$, $B$ thuộc đường tròn $(O)$ sao cho $\widehat{AOB}=110^\circ$. Số đo của cung lớn $AB$ là
	\choice
	{$110^\circ$}
	{$55^\circ$}
	{$125^\circ$}
	{\True $250^\circ$}
	\loigiai{
		Số đo cung nhỏ $AB$ bằng số đo góc ở tâm: $110^\circ$.\\
		Số đo cung lớn $AB$ là $360^\circ-110^\circ=250^\circ$.
	}
\end{ex}

% nguồn: dan_tu_AI (nhập 14/11/2026)
\begin{ex}%[9H5H14-2]%[ma:3c583acd]
	Cho đường tròn $(O;3\text{ cm})$ và ba điểm $A$, $B$, $C$ thuộc đường tròn sao cho tia $OB$ nằm giữa hai tia $OA$, $OC$, $\widehat{AOB}=80^\circ$ và $\widehat{BOC}=100^\circ$. Các khẳng định sau đúng hay sai?
	\choiceTF
	{\True Số đo cung nhỏ $AB$ bằng $80^\circ$}
	{Số đo cung lớn $AB$ bằng $260^\circ$}
	{\True $AC$ là một đường kính của đường tròn $(O)$}
	{\True Gọi $D$ là điểm đối xứng với $B$ qua $O$, khi đó số đo cung nhỏ $AD$ bằng $100^\circ$}
	\loigiai{
		\begin{itemchoice}
			\itemch Đúng. Số đo cung nhỏ $AB$ bằng góc ở tâm $\widehat{AOB}=80^\circ$.
			\itemch Sai. Số đo cung lớn $AB$ là $360^\circ-80^\circ=280^\circ$.
			\itemch Đúng. Vì tia $OB$ nằm giữa nên $\widehat{AOC}=80^\circ+100^\circ=180^\circ$, suy ra $A$, $O$, $C$ thẳng hàng và $AC$ là đường kính.
			\itemch Đúng. $D$ đối xứng với $B$ qua $O$ nên tia $OD$ là tia đối của tia $OB$, do đó $\widehat{AOD}=180^\circ-80^\circ=100^\circ$; số đo cung nhỏ $AD$ là $100^\circ$.
		\end{itemchoice}
	}
\end{ex}

% nguồn: dan_tu_AI (nhập 14/11/2026)
\begin{bt}%[9H5C14-2]%[ma:f817a95d]
	Mặt đồng hồ là một đường tròn có tâm $O$ tại trục quay của các kim. Lúc $3$ giờ $40$ phút, số đo góc ở tâm tạo bởi kim giờ và kim phút (góc không vượt quá $180^\circ$) là bao nhiêu độ?
	\shortans{130}
	\loigiai{
		Mỗi phút, kim phút quay $\dfrac{360^\circ}{60}=6^\circ$; mỗi giờ kim giờ quay $30^\circ$, mỗi phút quay $0{,}5^\circ$.\\
		Tính từ vị trí số $12$: kim phút quay được $40\cdot 6^\circ=240^\circ$; kim giờ quay được $3\cdot 30^\circ+40\cdot 0{,}5^\circ=110^\circ$.\\
		Góc ở tâm tạo bởi hai kim: $240^\circ-110^\circ=130^\circ$.
	}
\end{bt}

% nguồn: 2026-11-21_lop9_chuong5_bai14_17 (nhập 21/11/2026)
\begin{ex}%[9H5N14-2]%[ma:748a3afe]
% nguon: SGK Toán 9 KNTT tập 1 – Bài 14, Cung và dây của một đường tròn
	Trên đường tròn $(O)$, lấy hai điểm $A, B$ sao cho $\widehat{AOB} = 75^\circ$. Số đo cung lớn $\overset{\frown}{AB}$ bằng
	\choice
	{$75^\circ$}
	{$105^\circ$}
	{$150^\circ$}
	{\True $285^\circ$}
	\loigiai{
		Số đo cung nhỏ $\overset{\frown}{AB}$ bằng số đo góc ở tâm chắn cung đó là $75^\circ$.\\
		Suy ra số đo cung lớn $\overset{\frown}{AB}$ bằng $360^\circ - 75^\circ = 285^\circ$.
	}
\end{ex}

% nguồn: 2026-11-21_lop9_chuong5_bai14_17 (nhập 21/11/2026)
\begin{ex}%[9H5H14-2]%[ma:0ccb9b4d]
% nguon: SGK Toán 9 KNTT tập 1 – Bài 14, Cung và dây của một đường tròn
	Trên đường tròn $(O; 4\text{ cm})$, lấy hai điểm $A, B$ sao cho $\widehat{AOB} = 120^\circ$ và lấy điểm $C$ nằm trên cung nhỏ $\overset{\frown}{AB}$ sao cho số đo cung nhỏ $\overset{\frown}{AC}$ bằng $50^\circ$. Xét tính đúng, sai của các khẳng định sau:
	\choiceTF
	{\True Số đo cung nhỏ $\overset{\frown}{AB}$ bằng $120^\circ$}
	{Số đo cung lớn $\overset{\frown}{AB}$ bằng $220^\circ$}
	{\True Số đo cung nhỏ $\overset{\frown}{BC}$ bằng $70^\circ$}
	{\True Độ dài dây $AB$ bằng $4\sqrt{3}\text{ cm}$}
	\loigiai{
		\begin{itemchoice}
			\itemch Đúng. Số đo cung nhỏ $\overset{\frown}{AB}$ bằng $\widehat{AOB} = 120^\circ$.
			\itemch Sai. Số đo cung lớn $\overset{\frown}{AB}$ bằng $360^\circ - 120^\circ = 240^\circ$.
			\itemch Đúng. Vì $C$ nằm trên cung nhỏ $\overset{\frown}{AB}$ nên số đo cung nhỏ $\overset{\frown}{BC}$ bằng $120^\circ - 50^\circ = 70^\circ$.
			\itemch Đúng. Kẻ $OH \perp AB$ tại $H$, tam giác $OAB$ cân tại $O$ có $\widehat{AOH} = 60^\circ$ nên $AH = OA \cdot \sin 60^\circ = 4 \cdot \dfrac{\sqrt{3}}{2} = 2\sqrt{3}\text{ cm}$, suy ra $AB = 2AH = 4\sqrt{3}\text{ cm}$.
		\end{itemchoice}
	}
\end{ex}

% nguồn: 2026-11-21_lop9_chuong5_bai14_17 (nhập 21/11/2026)
\begin{bt}%[9H5V14-2]%[ma:bea46643]
% nguon: SGK Toán 9 KNTT tập 1 – Bài 14, Cung và dây của một đường tròn
	Lúc $4$ giờ đúng, kim giờ và kim phút của một đồng hồ tạo thành góc ở tâm chắn cung nhỏ trên mặt đồng hồ có số đo bằng bao nhiêu độ?
	\shortans{120}
	\loigiai{
		Mặt đồng hồ chia đường tròn thành $12$ cung bằng nhau, mỗi cung tương ứng với $1$ giờ có số đo là $\dfrac{360^\circ}{12} = 30^\circ$.\\
		Lúc $4$ giờ đúng, kim phút chỉ số $12$ và kim giờ chỉ số $4$, khoảng cách giữa hai kim là $4$ khoảng giờ nên số đo cung nhỏ bị chắn là $4 \cdot 30^\circ = 120^\circ$.
	}
\end{bt}

% nguồn: 2026-11-21_lop9_chuong5_bai14_15_bs (nhập 21/11/2026)
\begin{ex}%[9H5N14-2]%[ma:5eb2f184]
% nguon: SGK Toán 9 KNTT tập 1 – Bài 14, Cung và dây của một đường tròn
	Góc ở tâm chắn nửa đường tròn có số đo bằng
	\choice
	{$60^\circ$}
	{$90^\circ$}
	{\True $180^\circ$}
	{$360^\circ$}
	\loigiai{
		Nửa đường tròn có số đo bằng $180^\circ$, nên góc ở tâm chắn nửa đường tròn là góc bẹt có số đo $180^\circ$.
	}
\end{ex}

% nguồn: 2026-11-21_lop9_chuong5_bai14_15_bs (nhập 21/11/2026)
\begin{ex}%[9H5N14-2]%[ma:6ebcaf38]
% nguon: SGK Toán 9 KNTT tập 1 – Bài 14, Cung và dây của một đường tròn
	Trên đường tròn $(O)$, hai điểm $M, N$ chia đường tròn thành hai cung, trong đó cung nhỏ có số đo $110^\circ$. Số đo cung lớn $\overset{\frown}{MN}$ bằng
	\choice
	{$70^\circ$}
	{$110^\circ$}
	{$220^\circ$}
	{\True $250^\circ$}
	\loigiai{
		Số đo cung lớn $\overset{\frown}{MN}$ bằng $360^\circ - 110^\circ = 250^\circ$.
	}
\end{ex}

% nguồn: 2026-11-21_lop9_chuong5_bai14_15_bs (nhập 21/11/2026)
\begin{ex}%[9H5H14-2]%[ma:12ce71e3]
% nguon: SGK Toán 9 KNTT tập 1 – Bài 14, Cung và dây của một đường tròn
	Trên đường tròn $(O; 6\text{ cm})$, lấy hai điểm $A, B$ sao cho $\widehat{AOB} = 90^\circ$. Xét tính đúng, sai của các khẳng định sau:
	\choiceTF
	{\True Số đo cung nhỏ $\overset{\frown}{AB}$ bằng $90^\circ$}
	{\True Số đo cung lớn $\overset{\frown}{AB}$ bằng $270^\circ$}
	{Tam giác $OAB$ là tam giác đều}
	{\True Độ dài dây $AB$ bằng $6\sqrt{2}\text{ cm}$}
	\loigiai{
		\begin{itemchoice}
			\itemch Đúng. Số đo cung nhỏ $\overset{\frown}{AB}$ bằng số đo góc ở tâm $\widehat{AOB} = 90^\circ$.
			\itemch Đúng. Số đo cung lớn $\overset{\frown}{AB}$ bằng $360^\circ - 90^\circ = 270^\circ$.
			\itemch Sai. Vì $OA = OB = 6\text{ cm}$ và $\widehat{AOB} = 90^\circ$ nên tam giác $OAB$ là tam giác vuông cân tại $O$.
			\itemch Đúng. Theo định lí Pythagore trong tam giác vuông $OAB$: $AB = \sqrt{6^2 + 6^2} = 6\sqrt{2}\text{ cm}$.
		\end{itemchoice}
	}
\end{ex}

% nguồn: 2026-11-21_lop9_chuong5_bai14_15_bs (nhập 21/11/2026)
\begin{ex}%[9H5H14-2]%[ma:773e433c]
% nguon: SGK Toán 9 KNTT tập 1 – Bài 14, Cung và dây của một đường tròn
	Trên đường tròn $(O; 5\text{ cm})$, lấy ba điểm $A, B, C$ theo thứ tự đó sao cho số đo cung nhỏ $\overset{\frown}{AB}$ bằng $60^\circ$ và số đo cung nhỏ $\overset{\frown}{BC}$ bằng $60^\circ$. Xét tính đúng, sai của các khẳng định sau:
	\choiceTF
	{\True $\widehat{AOB} = 60^\circ$}
	{\True Số đo cung nhỏ $\overset{\frown}{AC}$ bằng $120^\circ$}
	{Độ dài dây $AB$ bằng $10\text{ cm}$}
	{\True Tứ giác $OABC$ là hình thoi có diện tích bằng $\dfrac{25\sqrt{3}}{2}\text{ cm}^2$}
	\loigiai{
		\begin{itemchoice}
			\itemch Đúng. Góc ở tâm $\widehat{AOB}$ bằng số đo cung nhỏ $\overset{\frown}{AB} = 60^\circ$.
			\itemch Đúng. Vì $B$ nằm trên cung nhỏ $\overset{\frown}{AC}$ nên số đo cung nhỏ $\overset{\frown}{AC}$ bằng $60^\circ + 60^\circ = 120^\circ$.
			\itemch Sai. Tam giác $OAB$ cân tại $O$ có $\widehat{AOB} = 60^\circ$ nên là tam giác đều, suy ra $AB = OA = 5\text{ cm}$.
			\itemch Đúng. Hai tam giác $OAB$ và $OBC$ đều cạnh $5\text{ cm}$ nên $OA = AB = BC = CO = 5\text{ cm}$, do đó $OABC$ là hình thoi có diện tích $2 \cdot \dfrac{5^2\sqrt{3}}{4} = \dfrac{25\sqrt{3}}{2}\text{ cm}^2$.
		\end{itemchoice}
	}
\end{ex}

% nguồn: 2026-11-21_lop9_chuong5_bai14_15_bs (nhập 21/11/2026)
\begin{bt}%[9H5V14-2]%[ma:d0afef49]
% nguon: SGK Toán 9 KNTT tập 1 – Bài 14, Cung và dây của một đường tròn
	Lúc $5$ giờ đúng, kim giờ và kim phút của một đồng hồ tạo thành góc ở tâm chắn cung nhỏ trên mặt đồng hồ có số đo bằng bao nhiêu độ?
	\shortans{150}
	\loigiai{
		Mỗi khoảng giờ trên mặt đồng hồ ứng với góc ở tâm $\dfrac{360^\circ}{12} = 30^\circ$.\\
		Lúc $5$ giờ đúng, kim phút chỉ số $12$ và kim giờ chỉ số $5$ nên góc ở tâm chắn cung nhỏ có số đo bằng $5 \cdot 30^\circ = 150^\circ$.
	}
\end{bt}

% nguồn: 2026-11-21_lop9_chuong5_bai14_15_bs (nhập 21/11/2026)
\begin{bt}%[9H5V14-2]%[ma:bf14304f]
% nguon: SGK Toán 9 KNTT tập 1 – Bài 14, Cung và dây của một đường tròn
	Trên đường tròn $(O)$, hai điểm $A, B$ chia đường tròn thành hai cung sao cho số đo cung lớn gấp $3$ lần số đo cung nhỏ. Tính số đo cung nhỏ $\overset{\frown}{AB}$ (đơn vị độ).
	\shortans{90}
	\loigiai{
		Gọi $x$ (độ) là số đo cung nhỏ $\overset{\frown}{AB}$ ($0 < x < 180$). Khi đó số đo cung lớn là $3x$.\\
		Ta có $x + 3x = 360 \Rightarrow 4x = 360 \Rightarrow x = 90$. Vậy số đo cung nhỏ $\overset{\frown}{AB}$ bằng $90^\circ$.
	}
\end{bt}
